\chapter{Symmetries and Reduction}
\label{ch:symmetries-reduction}

Classical lamphrodynamics features a hierarchy of symmetries acting on
the fields $(\Phi,\vec v,S)$: internal symmetries, spacetime
diffeomorphisms, and entropy--preserving reparametrizations.
At the derived level these are encoded as group actions on the plenum
stack.  Reduction of shifted symplectic structures along such actions
produces reduced phase spaces and determines the structure of the BV
complex.  In this chapter we formulate these ideas rigorously, develop
quotient stacks and reduction, and establish the geometric foundations
for BV--BFV boundary theory.



\section{Group Actions on the Plenum Stack}
\label{sec:group-actions}

Let $G$ be a (possibly infinite--dimensional) Lie group acting smoothly
on the configuration bundle of lamphrodynamic fields.
Passing to the derived stack level,
\[
G\curvearrowright\mathrm{Plenum}.
\]

\begin{definition}[Action on Derived Stacks]
A \emph{derived action} of $G$ on $\mathrm{Plenum}$
is a morphism of stacks
\[
\alpha\colon B G\times \mathrm{Plenum} \longrightarrow \mathrm{Plenum}
\]
compatible with group multiplication in the homotopical sense.
\end{definition}

\begin{remark}
The formalism of $BG$ allows the action to carry higher homotopies, so
derivations of the action may vanish only up to coherent homotopy.
\end{remark}



\section{Quotient Stacks and Derived Reduction}
\label{sec:quotients}

We recall the standard homotopy quotient construction.

\begin{definition}[Homotopy Quotient]
The \emph{derived quotient stack} is
\[
\mathrm{Plenum}\sslash G :=
\mathrm{Plenum}\times_{G} EG,
\]
where $EG\to BG$ is a contractible principal $G$--bundle.
\end{definition}

\begin{proposition}[Properties]
\label{prop:quotient-props}
The quotient stack $\mathrm{Plenum}\sslash G$
is a derived Artin stack with tangent complex of amplitude $[-1,1]$.
\end{proposition}

\begin{proof}[Proof (Sketch)]
The homotopy quotient preserves Artin representability and tangent
amplitude under mild hypotheses, using standard results in derived
geometry \cite{ToenVezzosi-HAGII,Calaque16}.
\end{proof}



\section{$(-1)$--Shifted Symplectic Reduction}
\label{sec:shifted-reduction}

Let $(\mathrm{Plenum},\omega)$ carry the $2$--shifted symplectic
structure constructed in Chapter~\ref{ch:plenum-construction}.
The quotient by a symmetry group $G$ inherits a shifted symplectic form.

\begin{theorem}[Shifted Symplectic Reduction]
\label{thm:shifted-reduction}
If $G$ acts on $\mathrm{Plenum}$ preserving $\omega$, then the
homotopy quotient $\mathrm{Plenum}\sslash G$ carries a
$(2-1)$--shifted symplectic structure, i.e.\ a $1$--shifted structure.
\end{theorem}

\begin{proof}[Proof (Sketch)]
This is a special case of shifted symplectic reduction
\cite{Calaque16}.  Preservation of $\omega$ ensures that
the descent to the quotient kills one degree of the shift.
\end{proof}

\begin{remark}
Under AKSZ descent on $\mathsf{T}[1]M$, the $1$--shift then produces a
$(-2)$--shift, which contributes to the ghost/antifield balance of
the BV theory.  Thus symmetry reduction influences the field content
after quantization.
\end{remark}



\section{Reduced Phase Spaces}
\label{sec:reduced-phase}

The classical (degree--$0$) phase space of lamphrodynamics is
\[
\mathcal{P}
:= \mathrm{Map}(M,\mathrm{Plenum})^{(0)},
\]
and the reduced phase space is obtained by quotienting by the derived
symmetry.

\begin{definition}[Reduced Classical Phase Space]
\[
\mathcal{P}_{\mathrm{red}}
:= \mathrm{Map}(M,\mathrm{Plenum}\sslash G)^{(0)}.
\]
\end{definition}

\begin{proposition}[Reduced Symplectic Structure]
\label{prop:reduced-symplectic}
$\mathcal{P}_{\mathrm{red}}$ inherits a classical symplectic form which
is the reduction of the classical form on $\mathcal{P}$.
\end{proposition}

\begin{proof}[Proof (Sketch)]
Restrict the shifted reduction of Theorem~\ref{thm:shifted-reduction}
to degree $0$.  Standard arguments in AKSZ theory apply.
\end{proof}



\section{Gauge Fixing and Derived Geometry}
\label{sec:gauge-fix}

Gauge fixing can be regarded as selecting a section of the projection
$\mathrm{Plenum}\to\mathrm{Plenum}\sslash G$.  Derived geometry provides
a canonical method of choosing such sections via homotopy transfer.

\begin{theorem}[Homotopy Gauge Fixing]
\label{thm:gauge-fix}
Under mild hypotheses on $\mathrm{Plenum}$, there exists a homotopy
splitting of the projection
\[
\mathrm{Plenum}\longrightarrow\mathrm{Plenum}\sslash G,
\]
so any gauge choice corresponds to a homotopy section.
\end{theorem}

\begin{proof}[Proof (Sketch)]
The projection is a homotopy epimorphism in the derived category; a
homotopy section exists by standard model category arguments.
\end{proof}

\begin{remark}
Homotopy gauge fixing means no physical content depends on gauge
choices; ghost/antifield sectors encode the homotopy redundancies
automatically.
\end{remark}



\section{BV--BFV Boundary Structure}
\label{sec:bv-bfv}

Boundary terms of the AKSZ action produce a BV--BFV (Batalin--Vilkovisky /
Batalin--Fradkin--Vilkovisky) theory on the boundary~$\partial M$,
encoding the symplectic structure of boundary degrees of freedom.

\begin{theorem}[BV--BFV Boundary Theorem]
\label{thm:bv-bfv}
If $M$ has nonempty boundary, the AKSZ--RSVP action induces a
$(-1)$--shifted symplectic form on the boundary field space and a
BV--BFV structure consistent with the reduced phase space of the
bulk theory.
\end{theorem}

\begin{proof}[Proof (Sketch)]
This follows from the general AKSZ--BV--BFV correspondence
\cite{CattaneoFelder2012}, applied to the $2$--shifted plenum and its
boundary reduction.
\end{proof}



\section{Summary and Forward Link}
\label{sec:summary13}

We have developed derived symmetries, quotient stacks, shifted
symplectic reduction, and reduced phase spaces for lamphrodynamic
fields.  Gauge fixing is homotopy--theoretic, and BV--BFV boundary
structures appear automatically.  These geometric foundations prepare
for the detailed classical comparisons of the next chapter.

